Bis zum Ziel bleibt die Strecke $s_\mathrm{Z}-s_\mathrm{M}$. Mit konstanter Geschwindigkeit braucht die Läuferin dafür
\begin{align}
 \Delta t &= \dtMF = \frac{\sZ-\sM}{\vM} \\
   &= \dtM \approx \dtMmP{0}{3}
\end{align}
also knapp drei Stunden. Zur Durchgangszeit $t_0 = \tMhO\ \tMmO$ (10:20 Uhr) addiert ergibt das $t_0+\Delta t = \tZ$ nach Mitternacht: Die Läuferin kommt um etwa $\result{\text{\hZI:\mZI~Uhr}}$ ins Ziel.
